\documentclass[11pt,letterpaper]{article}
\def\SuiteDocTitle{Source and Methodology Notes}
\def\SuiteDocCode{REF-01}
% Shared professional style: pdfLaTeX compatible, no external font files.
\usepackage[T1]{fontenc}
\usepackage[utf8]{inputenc}
\usepackage{lmodern}
\usepackage{microtype}
\usepackage[margin=0.83in,headheight=16pt,footskip=28pt]{geometry}
\usepackage[table]{xcolor}
\definecolor{navy}{HTML}{12304A}
\definecolor{teal}{HTML}{087F83}
\definecolor{gold}{HTML}{B78633}
\definecolor{blueback}{HTML}{F0F5F8}
\definecolor{tealback}{HTML}{EEF8F6}
\definecolor{amberback}{HTML}{FFF7E8}
\definecolor{grayback}{HTML}{F6F7F8}
\usepackage{array,tabularx,longtable,booktabs}
\usepackage{enumitem}
\usepackage{amssymb}
\usepackage{titlesec}
\usepackage{fancyhdr,lastpage}
\usepackage{tcolorbox}
\tcbuselibrary{breakable}
\usepackage{pgffor}
\usepackage{hyperref}
\hypersetup{colorlinks=true,linkcolor=teal,urlcolor=teal,pdfauthor={Engineering Learning Series},pdftitle={\SuiteDocTitle}}
\setlength{\parindent}{0pt}
\setlength{\parskip}{5.5pt}
\setlist{topsep=3pt}
\setcounter{tocdepth}{2}
\titleformat{\section}{\Large\bfseries\color{navy}}{\thesection}{0.7em}{}
\titleformat{\subsection}{\large\bfseries\color{teal}}{\thesubsection}{0.6em}{}
\titleformat{\subsubsection}{\normalsize\bfseries\color{navy}}{\thesubsubsection}{0.6em}{}
\pagestyle{fancy}
\fancyhf{}
\fancyhead[L]{\scriptsize\color{navy}AGILE REQUIREMENTS ENGINEERING}
\fancyhead[R]{\scriptsize\color{teal}\SuiteDocCode}
\fancyfoot[L]{\scriptsize Independent learning and implementation materials}
\fancyfoot[R]{\scriptsize Page \thepage\ of \pageref{LastPage}}
\renewcommand{\headrulewidth}{0.35pt}
\renewcommand{\footrulewidth}{0pt}
\newcommand{\SuiteHeader}[2]{%
{\color{teal}\small\bfseries ENGINEERING DOCUMENTATION SUITE \hfill #2}\par\vspace{8pt}
{\color{navy}\LARGE\bfseries\SuiteDocTitle}\par\vspace{5pt}
{\color{black!68}\normalsize #1}\par\vspace{8pt}
{\color{gold}\rule{\linewidth}{2pt}}\par\vspace{8pt}}
\newcommand{\WritingLines}[1]{\par\foreach \n in {1,...,#1}{\noindent\rule{\linewidth}{0.35pt}\par\vspace{0.28cm}}}
\newcommand{\FreeSpace}[1]{\begin{tcolorbox}[colback=white,colframe=black!25,height=#1,boxrule=0.5pt]\end{tcolorbox}}
\newcommand{\Field}[1]{\noindent\textbf{#1}\quad\rule{0.42\linewidth}{0.35pt}\par\vspace{0.3cm}}
\newcommand{\CheckItem}[1]{\noindent$\square$\quad #1\par\vspace{2pt}}
\newcommand{\SmallNote}[1]{{\footnotesize\color{black!62}#1}\par}
\newcommand{\Guidance}[1]{\begin{tcolorbox}[colback=blueback,colframe=teal,breakable,title=Guidance]#1\end{tcolorbox}}
\sloppy

\begin{document}
\SuiteHeader{Attribution | historical framing and original extensions}{Reference}
\tableofcontents\medskip\noindent\rule{\linewidth}{0.4pt}\medskip
\section{Primary supplied source}
Mike Cohn. User Stories Applied: For Agile Software Development. Addison-Wesley, 2004. ISBN 0-321-20568-5. User-provided PDF: 291 viewer pages. Printed book pages referenced in the crosswalk start after front matter. The material is copyrighted; this suite is original instruction and commentary, not a reproduction or replacement.

\section{Concept provenance}
\begin{longtable}{|p{0.20\linewidth}|p{0.29\linewidth}|p{0.37\linewidth}|}\hline
\textbf{Concept from the supplied book} & \textbf{Relevant chapters} & \textbf{Use in suite} \\ \hline\endfirsthead\hline
\textbf{Concept from the supplied book} & \textbf{Relevant chapters} & \textbf{Use in suite} \\ \hline\endhead
Card, Conversation, Confirmation; customer and developer collaboration & 1 & Technical reference and story card \\ \hline
INVEST and customer value & 2 & Story quality guide and scorecard \\ \hline
Role modeling, interviews, observation and proxies & 3–5 & Discovery workbook and Lab 01 \\ \hline
Acceptance tests specify intended behavior before coding & 6, 21 & Acceptance workbook and Lab 05 \\ \hline
Goal stories, vertical slicing, constraints and UI deferral & 7, 16 & Slicing template and nonfunctional constraints \\ \hline
Story points, release/iteration plans and velocity & 8–11, 19–20 & Planning workbook and Labs 06–07 \\ \hline
Contrasts with specifications, use cases and scenarios & 12–13 & Requirements reference \\ \hline
Story smells and Scrum adaptation & 14–15 & Scorecards and operating procedures \\ \hline
Extended sailing bookstore example & 17–21 & Methodological inspiration only; service-request examples are original \\ \hline
\end{longtable}
\section{What this source does not establish}
The 2004 book is not a source for current Scrum Guide rules, current ISO/IEC/IEEE standards, modern repository platform behavior, CI/CD product syntax, present-day security policy or empirical project benchmarks. All repository metadata patterns, Gherkin-style samples, quantitative teaching targets, governance gates and access-control demonstrations in this suite are independently constructed examples, not statements attributed to Cohn. Their suitability must be validated locally.

\section{Intended citation practice}
When using this material in internal reports, cite the book by author, title, year and relevant chapter, and cite original contemporary sources separately if a binding standard or organization-specific control is introduced. Do not cite the study guides as if they were the original text.


\end{document}
